\begin{bt}%[9D6H18-2]
% nguon: Tuyển tập đề tham khảo môn Toán Tuyển sinh 10 năm học 2025-2026 TP.HCM – Đề tham khảo Quận Bình Tân (Đề 1) – Câu 1
Cho hàm số $y=\dfrac{1}{4}x^2$ có đồ thị là parabol $(P)$.
\begin{enumerate}[a)]
	\item Vẽ đồ thị $(P)$.
	\item Tìm toạ độ những điểm $M$ thuộc $(P)$ có tung độ và hoành độ bằng nhau.
\end{enumerate}
\loigiai{
\begin{enumerate}[a)]
	\item Bảng giá trị tương ứng của $x$ và $y$:
	\begin{center}
	\begin{tabular}{|c|c|c|c|c|c|}
	\hline
	$x$ & $-4$ & $-2$ & $0$ & $2$ & $4$ \\ \hline
	$y=\dfrac{1}{4}x^2$ & $4$ & $1$ & $0$ & $1$ & $4$ \\ \hline
	\end{tabular}
	\end{center}
	Đồ thị $(P)$ là một đường parabol đi qua gốc toạ độ $O(0;0)$, nhận trục tung $Oy$ làm trục đối xứng và có bề lõm quay lên phía trên.
	\item Gọi $M(x_0; y_0)$ là điểm thuộc $(P)$.\\
	Theo giả thiết, điểm $M$ có tung độ và hoành độ bằng nhau nên $y_0 = x_0$.\\
	Do $M \in (P)$ nên ta có phương trình:
	$$ \dfrac{1}{4}x_0^2 = x_0 \iff x_0^2 - 4x_0 = 0 \iff x_0(x_0 - 4) = 0 \iff \left[\begin{aligned} x_0 &= 0 \\ x_0 &= 4. \end{aligned}\right. $$
	Với $x_0 = 0 \implies y_0 = 0$, ta được điểm $O(0;0)$.\\
	Với $x_0 = 4 \implies y_0 = 4$, ta được điểm $M(4;4)$.\\
	Vậy có hai điểm thoả mãn yêu cầu đề bài là $O(0;0)$ và $M(4;4)$.
\end{enumerate}
}
\end{bt}

\begin{bt}%[9D6H20-2]
% nguon: Tuyển tập đề tham khảo môn Toán Tuyển sinh 10 năm học 2025-2026 TP.HCM – Đề tham khảo Quận Bình Tân (Đề 1) – Câu 2
Cho phương trình bậc hai: $4x^2-2x-5=0$.
\begin{enumerate}[a)]
	\item Chứng tỏ phương trình luôn có hai nghiệm phân biệt $x_1, x_2$.
	\item Không giải phương trình, hãy tính giá trị của biểu thức:
	$$ B = \dfrac{x_1}{x_2-1} + \dfrac{x_2}{x_1-1}. $$
\end{enumerate}
\loigiai{
\begin{enumerate}[a)]
	\item Phương trình $4x^2-2x-5=0$ có các hệ số $a=4$, $b=-2$, $c=-5$.\\
	Vì $a \cdot c = 4 \cdot (-5) = -20 < 0$ nên phương trình luôn có hai nghiệm phân biệt $x_1, x_2$.
	\item Theo hệ thức Vi-ét, ta có:
	$$ \begin{cases} S = x_1+x_2 = -\dfrac{-2}{4} = \dfrac{1}{2} \\ P = x_1 x_2 = -\dfrac{5}{4}. \end{cases} $$
	Quy đồng mẫu thức biểu thức $B$:
	$$ B = \dfrac{x_1(x_1-1) + x_2(x_2-1)}{(x_2-1)(x_1-1)} = \dfrac{x_1^2 + x_2^2 - (x_1+x_2)}{x_1 x_2 - (x_1+x_2) + 1} = \dfrac{(x_1+x_2)^2 - 2x_1 x_2 - (x_1+x_2)}{x_1 x_2 - (x_1+x_2) + 1}. $$
	Thay $S = \dfrac{1}{2}$ và $P = -\dfrac{5}{4}$ vào biểu thức $B$, ta được:
	$$ B = \dfrac{\left(\dfrac{1}{2}\right)^2 - 2\left(-\dfrac{5}{4}\right) - \dfrac{1}{2}}{-\dfrac{5}{4} - \dfrac{1}{2} + 1} = \dfrac{\dfrac{1}{4} + \dfrac{5}{2} - \dfrac{1}{2}}{-\dfrac{3}{4}} = \dfrac{\dfrac{9}{4}}{-\dfrac{3}{4}} = -3. $$
\end{enumerate}
}
\end{bt}

\begin{bt}%[9D8H26-3]
% nguon: Tuyển tập đề tham khảo môn Toán Tuyển sinh 10 năm học 2025-2026 TP.HCM – Đề tham khảo Quận Bình Tân (Đề 1) – Câu 3
\begin{enumerate}[a)]
	\item Một chiếc hộp chứa $4$ tấm thẻ cùng loại được đánh số lần lượt là $1; 4; 7; 9$. Bạn Khuê và bạn Hương lần lượt mỗi người lấy ngẫu nhiên ra $1$ tấm thẻ từ hộp mà không trả lại. Tính xác suất của biến cố $A$: ``Tích các số ghi trên hai tấm thẻ lấy ra là một số lẻ''.
	\item Biểu đồ hình cột biểu diễn dữ liệu về hoạt động trong giờ ra chơi của học sinh lớp 9D: Đọc sách ($10$ học sinh), Ôn bài ($5$ học sinh), Đánh cầu lông ($4$ học sinh), Đá cầu ($8$ học sinh), Chơi cờ vua ($5$ học sinh), Nhảy dây ($2$ học sinh).\\
	Hãy phân tích dữ liệu từ biểu đồ trên để so sánh số học sinh tham gia hoạt động tại chỗ (đọc sách, ôn bài, chơi cờ vua) và hoạt động vận động (đánh cầu lông, đá cầu, nhảy dây) trong giờ ra chơi. Theo em, các bạn học sinh lớp 9D nên tăng cường loại hoạt động nào để có lợi hơn cho sức khoẻ?
\end{enumerate}
\loigiai{
\begin{enumerate}[a)]
	\item Trong $4$ tấm thẻ $\{1; 4; 7; 9\}$, có $3$ thẻ ghi số lẻ là $\{1; 7; 9\}$ và $1$ thẻ ghi số chẵn là $\{4\}$.\\
	Số phần tử của không gian mẫu (chọn lần lượt $2$ thẻ không hoàn lại) là:
	$$ n(\Omega) = 4 \cdot 3 = 12. $$
	Tích các số ghi trên hai tấm thẻ là số lẻ khi và chỉ khi cả hai thẻ lấy ra đều ghi số lẻ.\\
	Số kết quả thuận lợi cho biến cố $A$ là:
	$$ n(A) = 3 \cdot 2 = 6. $$
	Xác suất của biến cố $A$ là:
	$$ P(A) = \dfrac{n(A)}{n(\Omega)} = \dfrac{6}{12} = \dfrac{1}{2} = 0{,}5 \ (50\%). $$
	\item Từ biểu đồ cột, ta tính số lượng học sinh tham gia từng nhóm hoạt động:
	\begin{itemize}
		\item Số học sinh tham gia hoạt động tại chỗ là:
		$$ 10 + 5 + 5 = 20\text{ (học sinh)}. $$
		\item Số học sinh tham gia hoạt động vận động là:
		$$ 4 + 8 + 2 = 14\text{ (học sinh)}. $$
	\end{itemize}
	So sánh: Số học sinh lựa chọn hoạt động tại chỗ ($20$ bạn) nhiều hơn số học sinh tham gia hoạt động vận động ($14$ bạn).\\
	Ý kiến: Các bạn học sinh lớp 9D nên tăng cường các hoạt động vận động (như đá cầu, nhảy dây, đánh cầu lông) để giúp cơ thể giải toả căng thẳng sau các tiết học trên lớp, tăng cường tuần hoàn máu và nâng cao sức khoẻ thể chất.
\end{enumerate}
}
\end{bt}

\begin{bt}%[9H0H31-3]
% nguon: Tuyển tập đề tham khảo môn Toán Tuyển sinh 10 năm học 2025-2026 TP.HCM – Đề tham khảo Quận Bình Tân (Đề 1) – Câu 4
Một chiếc hộp có dạng hình lăng trụ đứng tam giác, đáy là một tam giác đều và chiều cao của lăng trụ bằng $9\text{ cm}$. Gọi $x\text{ (cm)}$ là độ dài cạnh đáy của chiếc hộp ($x > 0$).
\begin{enumerate}[a)]
	\item Viết biểu thức tính diện tích xung quanh của chiếc hộp theo $x$.
	\item Tính độ dài cạnh đáy của chiếc hộp biết diện tích xung quanh của hộp là $432\text{ cm}^2$.
\end{enumerate}
\loigiai{
\begin{enumerate}[a)]
	\item Đáy của lăng trụ đứng là tam giác đều cạnh $x\text{ (cm)}$, chu vi đáy là:
	$$ C_{\text{đáy}} = 3x\text{ (cm)}. $$
	Chiều cao lăng trụ là $h = 9\text{ cm}$. Diện tích xung quanh của chiếc hộp là:
	$$ S_{xq} = C_{\text{đáy}} \cdot h = 3x \cdot 9 = 27x\text{ (cm}^2\text{)}. $$
	\item Theo đề bài, diện tích xung quanh chiếc hộp là $432\text{ cm}^2$, ta có phương trình:
	$$ 27x = 432 \iff x = \dfrac{432}{27} = 16\text{ (cm)}. $$
	Vậy độ dài cạnh đáy của chiếc hộp là $16\text{ cm}$.
\end{enumerate}
}
\end{bt}

\begin{bt}%[9H0V31-3]
% nguon: Tuyển tập đề tham khảo môn Toán Tuyển sinh 10 năm học 2025-2026 TP.HCM – Đề tham khảo Quận Bình Tân (Đề 1) – Câu 5
Một que kem ốc quế gồm hai phần: phần kem bên trên có dạng hình cầu, phần ốc quế bên dưới có dạng hình nón. Biết rằng bán kính hình cầu và bán kính đáy của hình nón bằng nhau và đều bằng $2{,}5\text{ cm}$; chiều cao của hình nón gấp $3$ lần bán kính hình cầu. Tính thể tích của cả que kem đó (lấy $\pi = 3{,}14$ và làm tròn kết quả đến hàng đơn vị).
\loigiai{
Bán kính hình cầu và bán kính đáy hình nón là $R = 2{,}5\text{ cm}$.\\
Chiều cao của phần hình nón là:
$$ h = 3R = 3 \cdot 2{,}5 = 7{,}5\text{ cm}. $$
Thể tích của phần kem hình cầu là:
$$ V_{\text{cầu}} = \dfrac{4}{3}\pi R^3 = \dfrac{4}{3} \cdot 3{,}14 \cdot (2{,}5)^3 \approx 65{,}42\text{ cm}^3. $$
Thể tích của phần ốc quế hình nón là:
$$ V_{\text{nón}} = \dfrac{1}{3}\pi R^2 h = \dfrac{1}{3}\pi R^2 (3R) = \pi R^3 = 3{,}14 \cdot (2{,}5)^3 \approx 49{,}06\text{ cm}^3. $$
Tổng thể tích của que kem ốc quế là:
$$ V = V_{\text{cầu}} + V_{\text{nón}} = \dfrac{7}{3}\pi R^3 = \dfrac{7}{3} \cdot 3{,}14 \cdot (2{,}5)^3 \approx 114{,}48\text{ cm}^3. $$
Làm tròn kết quả đến hàng đơn vị, thể tích que kem là $114\text{ cm}^3$.
}
\end{bt}

\begin{bt}%[9D6V21-3]
% nguon: Tuyển tập đề tham khảo môn Toán Tuyển sinh 10 năm học 2025-2026 TP.HCM – Đề tham khảo Quận Bình Tân (Đề 1) – Câu 6
Một người dự định trồng $126$ cây xanh theo một kế hoạch thời gian định trước. Do thời tiết không thuận lợi nên thực tế mỗi ngày người đó trồng được ít hơn dự định $5$ cây, vì vậy thời gian hoàn thành công việc bị chậm mất $5$ ngày so với kế hoạch ban đầu. Hỏi trong thực tế mỗi ngày người đó đã trồng được bao nhiêu cây xanh?
\loigiai{
Gọi số cây người đó dự định trồng trong một ngày là $x$ ($x \in \mathbb{N}^*, x > 5$, cây/ngày).\\
Thời gian dự định trồng xong $126$ cây là:
$$ \dfrac{126}{x}\text{ (ngày)}. $$
Thực tế mỗi ngày người đó trồng được số cây là:
$$ x - 5\text{ (cây/ngày)}. $$
Thời gian thực tế để trồng xong $126$ cây là:
$$ \dfrac{126}{x - 5}\text{ (ngày)}. $$
Vì thực tế hoàn thành chậm mất $5$ ngày so với dự định nên ta có phương trình:
$$ \dfrac{126}{x - 5} - \dfrac{126}{x} = 5. $$
Quy đồng và khử mẫu thức ($x > 5$):
$$ 126x - 126(x - 5) = 5x(x - 5) \iff 630 = 5x^2 - 25x \iff 5x^2 - 25x - 630 = 0 \iff x^2 - 5x - 126 = 0. $$
Giải phương trình bậc hai:
$$ \Delta = (-5)^2 - 4 \cdot 1 \cdot (-126) = 25 + 504 = 529 = 23^2. $$
Phương trình có hai nghiệm:
$$ x_1 = \dfrac{5 + 23}{2} = 14\text{ (nhận)}, \quad x_2 = \dfrac{5 - 23}{2} = -9\text{ (loại)}. $$
Dự định mỗi ngày trồng $14$ cây. Vậy thực tế mỗi ngày người đó trồng được:
$$ 14 - 5 = 9\text{ (cây)}. $$
}
\end{bt}

\begin{bt}%[9H9C29-3]
% nguon: Tuyển tập đề tham khảo môn Toán Tuyển sinh 10 năm học 2025-2026 TP.HCM – Đề tham khảo Quận Bình Tân (Đề 1) – Câu 7
Cho đường tròn $(O;R)$ và một điểm $A$ nằm ngoài $(O)$ sao cho $OA = 2R$. Từ $A$ kẻ hai tiếp tuyến $AB, AC$ với đường tròn $(O)$ ($B, C$ là các tiếp điểm). Kẻ cát tuyến $ADE$ cắt đường tròn $(O)$ tại hai điểm phân biệt $D$ và $E$ ($D$ nằm giữa $A$ và $E$; tâm $O$ nằm trong góc $\widehat{BAE}$). Gọi $I$ là trung điểm của dây cung $DE$ ($OI \perp DE$ tại $I$).
\begin{enumerate}[a)]
	\item Chứng minh năm điểm $A, B, O, I, C$ cùng thuộc một đường tròn. Xác định tâm $K$ và bán kính của đường tròn này.
	\item Chứng minh $AB^2 = AD \cdot AE$.
	\item Tính diện tích phần hình tròn giới hạn bởi cung nhỏ $OC$ và dây cung $OC$ của đường tròn $(K)$ theo $R = 5{,}1\text{ cm}$ và lấy $\pi \approx 3{,}14$ (kết quả làm tròn đến chữ số thập phân thứ hai).
\end{enumerate}
\loigiai{
\begin{enumerate}[a)]
	\item Vì $AB, AC$ là hai tiếp tuyến của đường tròn $(O)$ nên $\widehat{ABO} = 90^\circ$ và $\widehat{ACO} = 90^\circ$.\\
	Vì $I$ là trung điểm của dây cung $DE$ không đi qua tâm nên $OI \perp DE$ tại $I \implies \widehat{AIO} = 90^\circ$.\\
	Ba điểm $B, I, C$ cùng nhìn đoạn thẳng $AO$ dưới một góc vuông, do đó năm điểm $A, B, O, I, C$ cùng thuộc đường tròn đường kính $AO$.\\
	Đường tròn này có tâm $K$ là trung điểm của đoạn thẳng $AO$, bán kính của đường tròn là:
	$$ R_K = \dfrac{AO}{2} = \dfrac{2R}{2} = R. $$
	\item Xét tam giác $ABD$ và tam giác $AEB$ có:
	$$ \widehat{BAD} \text{ là góc chung}, \quad \widehat{ABD} = \widehat{AEB} \text{ (góc tạo bởi tia tiếp tuyến và dây cung cùng chắn cung } BD). $$
	Suy ra $\Delta ABD \backsim \Delta AEB$ (g-g) $\implies \dfrac{AB}{AE} = \dfrac{AD}{AB} \implies AB^2 = AD \cdot AE$.
	\item Xét tam giác vuông $ACO$ vuông tại $C$, ta có $OC = R$ và $AO = 2R$.\\
	Suy ra $\sin \widehat{OAC} = \dfrac{OC}{AO} = \dfrac{R}{2R} = \dfrac{1}{2} \implies \widehat{OAC} = 30^\circ \implies \widehat{AOC} = 60^\circ$.\\
	Trong đường tròn $(K)$ đường kính $AO = 2R$, tâm $K$ là trung điểm $AO$ nên $KO = KC = R_K = R$.\\
	Mặt khác, tam giác $KOC$ có $KO = KC = R$ và $OC = R$, nên $\Delta KOC$ là một tam giác đều cạnh $R$.\\
	Do đó góc ở tâm của đường tròn $(K)$ chắn cung nhỏ $OC$ là $\widehat{OKC} = 60^\circ$.\\
	Diện tích hình quạt tròn chắn cung nhỏ $OC$ của đường tròn $(K)$ bán kính $R = 5{,}1\text{ cm}$ là:
	$$ S_{\text{quạt}} = \dfrac{\pi R^2 \cdot 60^\circ}{360^\circ} = \dfrac{3{,}14 \cdot (5{,}1)^2}{6} \approx 13{,}61\text{ cm}^2. $$
	Diện tích tam giác đều $KOC$ là:
	$$ S_{\Delta KOC} = \dfrac{R^2\sqrt{3}}{4} = \dfrac{(5{,}1)^2 \cdot \sqrt{3}}{4} \approx 11{,}26\text{ cm}^2. $$
	Diện tích phần hình tròn giới hạn bởi cung nhỏ $OC$ và dây cung $OC$ (hình viên phân) của đường tròn $(K)$ là:
	$$ S = S_{\text{quạt}} - S_{\Delta KOC} \approx 13{,}61 - 11{,}26 = 2{,}35\text{ cm}^2. $$
\end{enumerate}
}
\end{bt}

\begin{bt}%[9D6H18-2]
% nguon: Tuyển tập đề tham khảo môn Toán Tuyển sinh 10 năm học 2025-2026 TP.HCM – Đề tham khảo Quận Bình Tân (Đề 2) – Câu 1
Cho hàm số $y=-\dfrac{1}{2}x^2$ có đồ thị là parabol $(P)$.
\begin{enumerate}[a)]
	\item Vẽ đồ thị $(P)$ trên hệ trục toạ độ $Oxy$.
	\item Tìm toạ độ các điểm $A$ thuộc $(P)$ có tung độ bằng $-2$.
\end{enumerate}
\loigiai{
\begin{enumerate}[a)]
	\item Bảng giá trị tương ứng của $x$ và $y$:
	\begin{center}
	\begin{tabular}{|c|c|c|c|c|c|}
	\hline
	$x$ & $-4$ & $-2$ & $0$ & $2$ & $4$ \\ \hline
	$y=-\dfrac{1}{2}x^2$ & $-8$ & $-2$ & $0$ & $-2$ & $-8$ \\ \hline
	\end{tabular}
	\end{center}
	Đồ thị $(P)$ là một đường parabol đi qua gốc toạ độ $O(0;0)$, nhận trục tung $Oy$ làm trục đối xứng và có bề lõm quay xuống phía dưới.
	\item Gọi $A(x_0; y_0)$ là điểm thuộc $(P)$. Theo đề bài, tung độ của $A$ bằng $-2$ nên $y_0 = -2$.\\
	Do $A \in (P)$ nên ta có:
	$$ -\dfrac{1}{2}x_0^2 = -2 \iff x_0^2 = 4 \iff x_0 = \pm 2. $$
	Với $x_0 = 2 \implies A_1(2; -2)$.\\
	Với $x_0 = -2 \implies A_2(-2; -2)$.\\
	Vậy có hai điểm thoả mãn yêu cầu đề bài là $A_1(2; -2)$ và $A_2(-2; -2)$.
\end{enumerate}
}
\end{bt}

\begin{bt}%[9D6H20-2]
% nguon: Tuyển tập đề tham khảo môn Toán Tuyển sinh 10 năm học 2025-2026 TP.HCM – Đề tham khảo Quận Bình Tân (Đề 2) – Câu 2
Cho phương trình bậc hai: $4x^2-7x+2=0$.
\begin{enumerate}[a)]
	\item Chứng minh phương trình trên luôn có hai nghiệm phân biệt $x_1, x_2$.
	\item Không giải phương trình, hãy tính giá trị của biểu thức:
	$$ A = \left(\dfrac{2}{3}x_1 - x_2\right)\left(\dfrac{2}{3}x_1 + x_2\right) + \dfrac{5}{9}x_1^2 + 2x_2^2. $$
\end{enumerate}
\loigiai{
\begin{enumerate}[a)]
	\item Phương trình $4x^2-7x+2=0$ có các hệ số $a=4$, $b=-7$, $c=2$.\\
	Biệt thức $\Delta$ của phương trình:
	$$ \Delta = (-7)^2 - 4 \cdot 4 \cdot 2 = 49 - 32 = 17 > 0. $$
	Vì $\Delta > 0$ nên phương trình luôn có hai nghiệm phân biệt $x_1, x_2$.
	\item Theo hệ thức Vi-ét, ta có:
	$$ \begin{cases} S = x_1 + x_2 = \dfrac{7}{4} \\ P = x_1 x_2 = \dfrac{2}{4} = \dfrac{1}{2}. \end{cases} $$
	Khai triển và thu gọn biểu thức $A$:
	$$ A = \left(\dfrac{4}{9}x_1^2 - x_2^2\right) + \dfrac{5}{9}x_1^2 + 2x_2^2 = \left(\dfrac{4}{9} + \dfrac{5}{9}\right)x_1^2 + (-1 + 2)x_2^2 = x_1^2 + x_2^2. $$
	Biểu diễn theo tổng và tích các nghiệm:
	$$ A = (x_1 + x_2)^2 - 2x_1 x_2 = S^2 - 2P. $$
	Thay $S = \dfrac{7}{4}$ và $P = \dfrac{1}{2}$ vào biểu thức $A$:
	$$ A = \left(\dfrac{7}{4}\right)^2 - 2 \cdot \dfrac{1}{2} = \dfrac{49}{16} - 1 = \dfrac{33}{16} = 2{,}0625. $$
\end{enumerate}
}
\end{bt}

\begin{bt}%[9D6H21-3]
% nguon: Tuyển tập đề tham khảo môn Toán Tuyển sinh 10 năm học 2025-2026 TP.HCM – Đề tham khảo Quận Bình Tân (Đề 2) – Câu 3
Một khu đất hình chữ nhật có chiều rộng là $x\text{ (m)}$ ($x > 0$), chiều dài hơn chiều rộng $6\text{ m}$. Người ta xây một hồ nước hình vuông ở bên trong khu đất đó. Biết rằng phần diện tích đất còn lại của khu vườn sau khi xây hồ nước là $214\text{ m}^2$.
\begin{enumerate}[a)]
	\item Viết biểu thức $y$ biểu thị độ dài cạnh của hồ nước hình vuông theo $x$.
	\item Tính độ dài cạnh hình vuông của hồ nước khi biết $x = 14$ (kết quả làm tròn đến hàng phần mười của mét).
\end{enumerate}
\loigiai{
\begin{enumerate}[a)]
	\item Chiều rộng của khu đất là $x\text{ (m)}$, chiều dài của khu đất là $x+6\text{ (m)}$.\\
	Diện tích toàn bộ khu đất hình chữ nhật là:
	$$ S = x(x+6) = x^2 + 6x\text{ (m}^2\text{)}. $$
	Gọi $y$ là độ dài cạnh của hồ nước hình vuông ($y > 0$, đơn vị: mét). Diện tích hồ nước là $y^2\text{ (m}^2\text{)}$.\\
	Vì phần diện tích còn lại là $214\text{ m}^2$ nên ta có:
	$$ x^2 + 6x - y^2 = 214 \iff y^2 = x^2 + 6x - 214. $$
	Vì $y > 0$ nên biểu thức tính độ dài cạnh hồ nước theo $x$ là:
	$$ y = \sqrt{x^2 + 6x - 214}\text{ (m)}. $$
	\item Thay $x = 14$ vào biểu thức của $y$, ta được:
	$$ y = \sqrt{14^2 + 6 \cdot 14 - 214} = \sqrt{196 + 84 - 214} = \sqrt{66} \approx 8{,}1\text{ (m)}. $$
	Vậy độ dài cạnh của hồ nước hình vuông là khoảng $8{,}1\text{ m}$.
\end{enumerate}
}
\end{bt}

\begin{bt}%[9H0V31-3]
% nguon: Tuyển tập đề tham khảo môn Toán Tuyển sinh 10 năm học 2025-2026 TP.HCM – Đề tham khảo Quận Bình Tân (Đề 2) – Câu 4
Một chiếc bình bằng thuỷ tinh hình trụ có đường kính đáy là $40\text{ cm}$ và chiều cao của bình là $60\text{ cm}$.
\begin{enumerate}[a)]
	\item Tính thể tích nước cần đổ vào bình ban đầu để mực nước trong bình cao $50\text{ cm}$ (biết lúc đầu bình chưa có nước).
	\item Người ta thả một số viên bi hình cầu bằng kim loại (không thấm nước) có cùng bán kính là $6\text{ cm}$ vào bình. Hỏi có thể thả vào bình nhiều nhất bao nhiêu viên bi như thế để nước trong bình không bị tràn ra ngoài? (Biết các viên bi khi thả vào đều chìm hoàn toàn xuống đáy bình, lấy $\pi = 3{,}14$).
\end{enumerate}
\loigiai{
\begin{enumerate}[a)]
	\item Bán kính đáy của chiếc bình hình trụ là $R = \dfrac{40}{2} = 20\text{ cm}$.\\
	Thể tích nước đổ vào để mực nước cao $h_1 = 50\text{ cm}$ là:
	$$ V_{\text{nước}} = \pi R^2 h_1 = 3{,}14 \cdot 20^2 \cdot 50 = 62\,800\text{ (cm}^3\text{)}. $$
	\item Thể tích toàn phần của bình hình trụ là:
	$$ V_{\text{bình}} = \pi R^2 h = 3{,}14 \cdot 20^2 \cdot 60 = 75\,360\text{ (cm}^3\text{)}. $$
	Phần thể tích còn trống phía trên mặt nước của bình là:
	$$ V_{\text{trống}} = V_{\text{bình}} - V_{\text{nước}} = 75\,360 - 62\,800 = 12\,560\text{ (cm}^3\text{)}. $$
	Thể tích của một viên bi hình cầu bán kính $r = 6\text{ cm}$ là:
	$$ V_{\text{bi}} = \dfrac{4}{3}\pi r^3 = \dfrac{4}{3} \cdot 3{,}14 \cdot 6^3 = 904{,}32\text{ (cm}^3\text{)}. $$
	Số viên bi tối đa có thể thả vào để nước không bị tràn ra ngoài là:
	$$ N \le \dfrac{V_{\text{trống}}}{V_{\text{bi}}} = \dfrac{12\,560}{904{,}32} \approx 13{,}88. $$
	Vì số viên bi phải là số nguyên dương nên có thể thả vào bình nhiều nhất là $13$ viên bi.
\end{enumerate}
}
\end{bt}

\begin{bt}%[9D1H3-4]
% nguon: Tuyển tập đề tham khảo môn Toán Tuyển sinh 10 năm học 2025-2026 TP.HCM – Đề tham khảo Quận Bình Tân (Đề 2) – Câu 5
Trong đợt khuyến mãi nhân dịp kỷ niệm ngày Quốc khánh 2/9, siêu thị A giảm giá $20\%$ cho một thùng nước ngọt và giảm $15\%$ cho một thùng sữa tươi so với giá niêm yết. Một khách hàng mua $2$ thùng nước ngọt và $1$ thùng sữa tươi phải trả tổng số tiền là $362\,000$ đồng. Trong khung giờ vàng cùng ngày, siêu thị giảm sâu hơn: một thùng nước ngọt được giảm giá $30\%$ còn một thùng sữa tươi được giảm giá $25\%$ so với giá niêm yết. Một khách hàng khác đã mua $3$ thùng nước ngọt và $2$ thùng sữa tươi trong khung giờ vàng chỉ phải trả số tiền là $552\,000$ đồng. Tính giá niêm yết ban đầu của mỗi thùng nước ngọt và mỗi thùng sữa tươi.
\loigiai{
Gọi $x$ và $y$ lần lượt là giá niêm yết ban đầu của một thùng nước ngọt và một thùng sữa tươi ($x, y > 0$, đơn vị: đồng).\\
Trong đợt khuyến mãi thông thường:
\begin{itemize}
	\item Giá một thùng nước ngọt sau giảm $20\%$ là: $(1 - 0{,}2)x = 0{,}8x$ (đồng).
	\item Giá một thùng sữa tươi sau giảm $15\%$ là: $(1 - 0{,}15)y = 0{,}85y$ (đồng).
\end{itemize}
Khách hàng mua $2$ thùng nước ngọt và $1$ thùng sữa tươi trả $362\,000$ đồng, ta có phương trình:
$$ 2 \cdot 0{,}8x + 0{,}85y = 362\,000 \iff 1{,}6x + 0{,}85y = 362\,000. $$
Trong khung giờ vàng:
\begin{itemize}
	\item Giá một thùng nước ngọt sau giảm $30\%$ là: $(1 - 0{,}3)x = 0{,}7x$ (đồng).
	\item Giá một thùng sữa tươi sau giảm $25\%$ là: $(1 - 0{,}25)y = 0{,}75y$ (đồng).
\end{itemize}
Khách hàng mua $3$ thùng nước ngọt và $2$ thùng sữa tươi trả $552\,000$ đồng, ta có phương trình:
$$ 3 \cdot 0{,}7x + 2 \cdot 0{,}75y = 552\,000 \iff 2{,}1x + 1{,}5y = 552\,000. $$
Ta có hệ phương trình bậc nhất hai ẩn:
$$ \begin{cases} 1{,}6x + 0{,}85y = 362\,000 \\ 2{,}1x + 1{,}5y = 552\,000 \end{cases} \iff \begin{cases} 32x + 17y = 7\,240\,000 \\ 7x + 5y = 1\,840\,000 \end{cases} $$
Nhân phương trình thứ hai với $3{,}4$ rồi trừ vế theo vế:
$$ \begin{cases} 32x + 17y = 7\,240\,000 \\ 23{,}8x + 17y = 6\,256\,000 \end{cases} \iff \begin{cases} 8{,}2x = 984\,000 \\ 7x + 5y = 1\,840\,000 \end{cases} \iff \begin{cases} x = 120\,000\text{ (thoả mãn)} \\ y = 200\,000\text{ (thoả mãn)}. \end{cases} $$
Vậy giá niêm yết ban đầu của một thùng nước ngọt là $120\,000$ đồng và của một thùng sữa tươi là $200\,000$ đồng.
}
\end{bt}

\begin{bt}%[9D8H26-3]
% nguon: Tuyển tập đề tham khảo môn Toán Tuyển sinh 10 năm học 2025-2026 TP.HCM – Đề tham khảo Quận Bình Tân (Đề 2) – Câu 6
Trong một nhóm gồm $10$ học sinh lớp 9 có $5$ bạn học trường THCS Quang Trung, $3$ bạn học trường THCS Nguyễn Huệ và $2$ bạn học trường THCS Tây Sơn. Người ta chọn ngẫu nhiên $1$ bạn học sinh trong nhóm $10$ bạn đó.
\begin{enumerate}[a)]
	\item Không gian mẫu của phép thử có bao nhiêu phần tử?
	\item Tính xác suất của các biến cố sau:
	\begin{itemize}
		\item Biến cố $A$: ``Bạn học sinh được chọn học trường THCS Quang Trung''.
		\item Biến cố $B$: ``Bạn học sinh được chọn không học trường THCS Tây Sơn''.
	\end{itemize}
\end{enumerate}
\loigiai{
\begin{enumerate}[a)]
	\item Phép thử là chọn ngẫu nhiên $1$ bạn học sinh từ nhóm $10$ học sinh, nên số phần tử của không gian mẫu là:
	$$ n(\Omega) = 10. $$
	\item Tính xác suất của từng biến cố:
	\begin{itemize}
		\item Có $5$ bạn học sinh trường THCS Quang Trung, do đó số kết quả thuận lợi cho biến cố $A$ là $n(A) = 5$.\\
		Xác suất của biến cố $A$ là:
		$$ P(A) = \dfrac{n(A)}{n(\Omega)} = \dfrac{5}{10} = \dfrac{1}{2} = 0{,}5 \ (50\%). $$
		\item Biến cố $B$ là bạn học sinh được chọn không học trường Tây Sơn, tức là bạn đó học trường Quang Trung hoặc Nguyễn Huệ.\\
		Số kết quả thuận lợi cho biến cố $B$ là:
		$$ n(B) = 10 - 2 = 8. $$
		Xác suất của biến cố $B$ là:
		$$ P(B) = \dfrac{n(B)}{n(\Omega)} = \dfrac{8}{10} = \dfrac{4}{5} = 0{,}8 \ (80\%). $$
	\end{itemize}
\end{enumerate}
}
\end{bt}

\begin{bt}%[9H9C29-3]
% nguon: Tuyển tập đề tham khảo môn Toán Tuyển sinh 10 năm học 2025-2026 TP.HCM – Đề tham khảo Quận Bình Tân (Đề 2) – Câu 7
Từ điểm $M$ nằm bên ngoài đường tròn $(O;R)$, vẽ hai tiếp tuyến $MA, MC$ với đường tròn $(O)$ ($A, C$ là các tiếp điểm). Vẽ đường kính $AB$ của $(O)$. Đoạn thẳng $OM$ cắt $AC$ tại điểm $I$.
\begin{enumerate}[a)]
	\item Chứng minh bốn điểm $A, M, C, O$ cùng thuộc một đường tròn và $MO \perp AC$ tại $I$.
	\item Kẻ dây cung $CD$ vuông góc với $AB$ tại $D$ ($D \in AB$). Đoạn thẳng $BM$ cắt đường tròn $(O)$ tại $N$ ($N \ne B$). Chứng minh hai tam giác $\Delta MAO$ và $\Delta CDB$ đồng dạng, và $MC^2 = MI \cdot MO$.
	\item Gọi $K$ là giao điểm của $BM$ và $CD$. Khi $OM = AB = 2R$, hãy tính theo $R$ diện tích tam giác $MIK$.
\end{enumerate}
\loigiai{
\begin{enumerate}[a)]
	\item Vì $MA, MC$ là hai tiếp tuyến của $(O)$ nên $\widehat{MAO} = 90^\circ$ và $\widehat{MCO} = 90^\circ$.\\
	Tứ giác $MAOC$ có $\widehat{MAO} + \widehat{MCO} = 90^\circ + 90^\circ = 180^\circ$, suy ra bốn điểm $A, M, C, O$ cùng thuộc đường tròn đường kính $MO$.\\
	Theo tính chất của hai tiếp tuyến cắt nhau, ta có $MA = MC$ và $OA = OC = R$, nên $MO$ là đường trung trực của đoạn thẳng $AC \implies MO \perp AC$ tại trung điểm $I$ của $AC$.
	\item Tam giác $ACB$ nội tiếp đường tròn $(O)$ đường kính $AB$ nên $\widehat{ACB} = 90^\circ$.\\
	Trong tam giác vuông $ACB$ vuông tại $C$ có đường cao $CD \perp AB$, ta có $\widehat{DCB} = \widehat{CAB}$ (cùng phụ với $\widehat{ACD}$).\\
	Mặt khác, $MA \perp AB$ và $CD \perp AB$ nên $MA \parallel CD \implies \widehat{AOM} = \widehat{CAB}$ (do cùng phụ với góc $\widehat{OAC}$, hoặc $\Delta MAO$ vuông tại $A$ có $\widehat{AMO} = \widehat{OAC}$).\\
	Xét hai tam giác vuông $\Delta MAO$ (vuông tại $A$) và $\Delta CDB$ (vuông tại $D$):
	$$ \widehat{MAO} = \widehat{CDB} = 90^\circ, \quad \widehat{AOM} = \widehat{DCB}. $$
	Suy ra $\Delta MAO \backsim \Delta CDB$ (g-g).\\
	Trong tam giác vuông $MCO$ vuông tại $C$ có đường cao $CI \perp MO$, áp dụng hệ thức lượng ta có:
	$$ MC^2 = MI \cdot MO. $$
	\item Khi $OM = AB = 2R$, trong tam giác vuông $MAO$ vuông tại $A$ có $OA = R, OM = 2R$:
	$$ MA = \sqrt{OM^2 - OA^2} = \sqrt{(2R)^2 - R^2} = R\sqrt{3}. $$
	Suy ra $\sin \widehat{AMO} = \dfrac{OA}{OM} = \dfrac{R}{2R} = \dfrac{1}{2} \implies \widehat{AMO} = 30^\circ \implies \widehat{AOM} = 60^\circ$.\\
	Vì $CD \parallel MA$ (cùng vuông góc với $AB$), áp dụng định lý Ta-lét trong tam giác $ABM$ có $KD \parallel MA$:
	$$ \dfrac{KD}{MA} = \dfrac{BD}{AB}. $$
	Vì $\widehat{AOM} = 60^\circ$ nên $\widehat{AOC} = 2\widehat{AOM} = 120^\circ \implies \widehat{BOC} = 180^\circ - 120^\circ = 60^\circ$.\\
	Trong tam giác vuông $ODC$ tại $D$:
	$$ OD = OC \cdot \cos 60^\circ = \dfrac{R}{2} \implies BD = OB - OD = R - \dfrac{R}{2} = \dfrac{R}{2}. $$
	Do đó:
	$$ \dfrac{BD}{AB} = \dfrac{\dfrac{R}{2}}{2R} = \dfrac{1}{4} \implies KD = \dfrac{1}{4}MA = \dfrac{R\sqrt{3}}{4}. $$
	Chiều dài $CD = OC \cdot \sin 60^\circ = \dfrac{R\sqrt{3}}{2}$, suy ra $CK = CD - KD = \dfrac{R\sqrt{3}}{2} - \dfrac{R\sqrt{3}}{4} = \dfrac{R\sqrt{3}}{4}$.\\
	Do đó $K$ là trung điểm của $CD$.\\
	Mặt khác, $I$ là trung điểm của $AC$, do đó $IK$ là đường trung bình của tam giác $ACD$, suy ra:
	$$ IK \parallel AD \parallel AB \quad \text{và} \quad IK = \dfrac{1}{2}AD = \dfrac{1}{2}(OA + OD) = \dfrac{1}{2}\left(R + \dfrac{R}{2}\right) = \dfrac{3R}{4}. $$
	Vì $IK \parallel AB$ mà $MA \perp AB$ nên khoảng cách từ đỉnh $M$ đến đường thẳng $IK$ là:
	$$ h = d(M, IK) = MA - d(IK, AB) = MA - KD = R\sqrt{3} - \dfrac{R\sqrt{3}}{4} = \dfrac{3R\sqrt{3}}{4}. $$
	Vậy diện tích của tam giác $MIK$ là:
	$$ S_{\Delta MIK} = \dfrac{1}{2} \cdot IK \cdot h = \dfrac{1}{2} \cdot \dfrac{3R}{4} \cdot \dfrac{3R\sqrt{3}}{4} = \dfrac{9\sqrt{3}}{32}R^2. $$
\end{enumerate}
}
\end{bt}
